% Pandoc template for the Quarto resume.
%
% This replaces Quarto's stock LaTeX template entirely. That is deliberate: the stock
% template loads geometry and hyperref with its own options, and resume.cls already
% loads both with ours, which raises an option clash. Everything typographic lives in
% resume.cls -- this file only bridges pandoc's output into it.

\documentclass{resume}

% Commands pandoc emits that the stock template would otherwise define.
\providecommand{\tightlist}{%
  \setlength{\itemsep}{0pt}\setlength{\parskip}{0pt}}
\providecommand{\passthrough}[1]{#1}
\providecommand{\pandocbounded}[1]{#1}

\hypersetup{pdftitle={Gage Rowden's Resume},pdfauthor={Gage Rowden}}

\makeatletter
\@ifpackageloaded{caption}{}{\usepackage{caption}}
\AtBeginDocument{%
\ifdefined\contentsname
  \renewcommand*\contentsname{Table of contents}
\else
  \newcommand\contentsname{Table of contents}
\fi
\ifdefined\listfigurename
  \renewcommand*\listfigurename{List of Figures}
\else
  \newcommand\listfigurename{List of Figures}
\fi
\ifdefined\listtablename
  \renewcommand*\listtablename{List of Tables}
\else
  \newcommand\listtablename{List of Tables}
\fi
\ifdefined\figurename
  \renewcommand*\figurename{Figure}
\else
  \newcommand\figurename{Figure}
\fi
\ifdefined\tablename
  \renewcommand*\tablename{Table}
\else
  \newcommand\tablename{Table}
\fi
}
\@ifpackageloaded{float}{}{\usepackage{float}}
\floatstyle{ruled}
\@ifundefined{c@chapter}{\newfloat{codelisting}{h}{lop}}{\newfloat{codelisting}{h}{lop}[chapter]}
\floatname{codelisting}{Listing}
\newcommand*\listoflistings{\listof{codelisting}{List of Listings}}
\makeatother
\makeatletter
\makeatother
\makeatletter
\@ifpackageloaded{caption}{}{\usepackage{caption}}
\@ifpackageloaded{subcaption}{}{\usepackage{subcaption}}
\makeatother

\begin{document}


\begin{header}

    \fontsize{25 pt}{25 pt}\selectfont Gage Rowden \textbf{Applied Data Scientist}\\
    \vspace{5 pt}
    \normalsize
    \mbox{\hrefWithoutArrow{mailto:gage.rowden1145@gmail.com}{\faEnvelope{ }gage.rowden1145@gmail.com}}
    \AND
    \mbox{\hrefWithoutArrow{tel:+1-806-577-8008}{\faPhone{ }+1-806-577-8008}}
    \AND
    \mbox{\hrefWithoutArrow{https://github.com/gage1145}{\faGithub{ }gage1145}}
    \AND
    \mbox{\hrefWithoutArrow{https://orcid.org/0000-0002-7517-0480}{\faOrcid{ }ORCID}}
    \AND
    \mbox{\hrefWithoutArrow{https://www.linkedin.com/in/gagerowden}{\faLinkedin{ }gagerowden}}
\end{header}

\vspace{5 pt - 0.3 cm}

\section{Professional Summary}\label{professional-summary}

\begin{centering}
\Large

Applied Data Scientist specializing in high-throughput biological data
and production analytics systems with extensive experience designing
end-to-end data workflows, developing custom R packages, and building
cloud-backed applications. Skilled in Python, SQL, AWS, and statistical
modeling, with a strong foundation in experimental design and
high-throughput data.

\end{centering}

\section{Experience}\label{experience}

\begin{twocolentry}{Jan 2024--present}

\textbf{Lead Data Scientist}, Priogen Corporation --- St.~Paul, MN

\end{twocolentry}

\begin{onecolentry}

\begin{highlights}
    \item Architected and deployed end-to-end \textbf{data analysis pipelines} to
process high-throughput diagnostic datasets (up to 100 samples/week,
$\sim$240k rows of time-series data/project), reducing analysis time
from $\sim$4 hours to seconds.
    \item Led \textbf{full-stack} development of a \textbf{Python}-based
\textbf{LIMS} (PriLIMS) using \textbf{Django}, automating client data
intake, tracking, and reporting, eliminating 10 hours/week of manual
administrative work.
    \item Implemented cloud-based infrastructure on \textbf{AWS} to support
scalable \textbf{SQL} storage and analysis, improving system reliability
and enabling secure multi-user access.
    \item Introduced \textbf{Git}-based version control and \textbf{CI} workflows
across projects, improving collaboration efficiency and significantly
reducing deployment errors.
\end{highlights}

\end{onecolentry}

\begin{twocolentry}{Aug 2022--present}

\textbf{Researcher IV}, MNPRO, University of Minnesota --- St.~Paul, MN

\end{twocolentry}

\begin{onecolentry}

\begin{highlights}
    \item Developed a custom \textbf{R} package and automation scripts to
standardize RT-QuIC data analysis, reducing hours of analysis time to
seconds and increasing reproducibility across collaborating labs.
    \item Streamlined \textbf{data processing} and reporting workflows using
\textbf{R} and \textbf{Python} tools, reducing turnaround time for
experimental results by \textbf{99\%}.
    \item Built and evaluated \textbf{statistical models} to classify diagnostic
outcomes, improving disease detection confidence.
    \item Standardized cross-lab analytics and built reusable analytical
frameworks adopted by multiple institutions.
    \item Led laboratory establishment, expansion, and workflow redesign,
increasing research capacity.
\end{highlights}

\end{onecolentry}

\section{Skills}\label{skills}

\begin{skills}
    \item Python
    \item Pandas \& NumPy
    \item scikit-learn
    \item R Tidyverse
    \item SQL
    \item ETL pipeline development
    \item Time-series analysis
    \item Statistical analysis
    \item Experimental design
    \item Data visualization
    \item AWS (cloud infrastructure)
    \item Django \& REST API
    \item Docker
    \item Git \& GitHub
    \item GitHub Actions (CI/CD)
\end{skills}

\section{Projects}\label{projects}

\begin{software}{\href{https://github.com/gage1145/quicR}{github.com/gage1145/quicR}}

\textbf{quicR: R Library for Streamlined Data Handling of RT-QuIC
Assays}

\end{software}

\begin{highlights}
    \item Developed an R package for the extraction, manipulation, and analysis of
RT-QuIC data.
    \item Stack: R, Tidyverse, RTools, Git, GitHub
\end{highlights}

\begin{software}{\href{https://priogen.xyz}{priogen.xyz}}

\textbf{PriLIMS: A Priogen Laboratory Information Management System}

\end{software}

\begin{highlights}
    \item Designed and implemented a Django framework LIMS for Priogen Corp.
    \item Stack: Python, R, Django, SQL, AWS, Git, GitHub, REST API, HTML, CSS,
Javascript
\end{highlights}

\begin{software}{\href{https:://modeling.gagerowden.com}{modeling.gagerowden.com}}

\textbf{Non-linear Kinetic Modeling of Protein Mis-Folding Assays}

\end{software}

\begin{highlights}
    \item Built a non-linear least-squares model that accurately predicts the
kinetic behavior of protein mis-folding assays.
    \item Stack: R, Tidyverse, Quarto, Observable JS, Git, GitHub
\end{highlights}


\end{document}
